\cleardoublepage
\chapter{PELAKSANAAN PEKERJAAN}
\label{chap:pekerjaan}

\section{Gambaran Umum Pekerjaan}
Pekerjaan penulis selama Kerja Praktik dapat dikelompokkan menjadi tiga deliverable utama dan satu kelompok pembelajaran pendukung. Ringkasan tersebut ditunjukkan pada Tabel \ref{tab:ringkasan-pekerjaan}.

\begin{table}[htbp]
\caption{Ringkasan pekerjaan selama Kerja Praktik}
\label{tab:ringkasan-pekerjaan}
\centering
\begin{tabularx}{\textwidth}{|p{0.23\textwidth}|p{0.32\textwidth}|X|}
\hline
\textbf{Pekerjaan} & \textbf{Fokus} & \textbf{Output} \\
\hline
Query dan ETL campaign & Menerjemahkan kebutuhan bisnis menjadi proses data & Query yang lebih terstruktur, validasi hasil, dan job ETL yang diuji untuk proses berkala \\
\hline
Analisis campaign financing & Membaca efektivitas campaign pada beberapa tingkat agregasi & Ringkasan deskriptif, perbandingan performa, dan bahan rekomendasi \\
\hline
Spatial clustering log aplikasi & Membaca pola lokasi pada data log aplikasi digital & Pipeline clustering, hasil agregat lokasi, reverse geocoding centroid, dan visualisasi \\
\hline
Pembelajaran pendukung & Memahami alur data, cleansing, leads, dan konteks bisnis & Catatan kerja, validasi konsep, dan pemahaman hubungan antara data dan proses campaign \\
\hline
\end{tabularx}
\end{table}

Ketiga pekerjaan utama memiliki bentuk yang berbeda, tetapi proses berpikirnya mirip. Penulis harus memahami pertanyaan bisnis, menentukan data yang relevan, memeriksa kualitas data, memilih cara pengolahan, lalu menyajikan hasil dengan batas interpretasi yang jelas.

\section{Pemahaman Awal terhadap Goal, Data, dan Kebutuhan Bisnis}
Pada awal Kerja Praktik, penulis mempelajari istilah bisnis, produk, campaign, leads, serta hubungan antar lingkungan data. Catatan kerja juga digunakan untuk memahami perbedaan fungsi Data Warehouse, Enterprise Data Lake, dan Data Mart pada konteks pekerjaan yang ditemui.

Pembelajaran yang paling penting pada tahap ini adalah menentukan goal sebelum mulai menarik data. Penulis perlu memahami keputusan atau output apa yang ingin dihasilkan. Dari goal tersebut, penulis menurunkan informasi yang dibutuhkan, metrik atau kriteria yang perlu dihitung, entitas yang terlibat, lalu atribut yang harus tersedia.

Setelah kebutuhan informasi jelas, penulis memetakan atribut tersebut ke kolom dan sumber tabel yang relevan. Tahap berikutnya adalah menentukan relasi antartabel, grain data, periode waktu, kondisi filter, flag, penanganan nilai kosong, serta kemungkinan duplikasi akibat \textit{join}. Jika satu atribut belum tersedia atau definisinya belum jelas, penulis perlu kembali ke kebutuhan bisnis untuk melakukan klarifikasi.

Urutan tersebut membuat proses query tidak dimulai dari tabel yang tersedia, tetapi dari goal yang ingin dicapai dan data yang dibutuhkan untuk mencapainya. Query, transformasi, dan analisis kemudian menjadi implementasi dari kebutuhan yang sudah dipetakan. Alur penurunan kebutuhan tersebut diringkas pada Gambar \ref{fig:goal-data}.

\begin{figure}[htbp]
\centering
\small
\setlength{\fboxsep}{6pt}
\fbox{\parbox{0.78\textwidth}{\centering \textbf{Goal bisnis}\\Keputusan atau output yang ingin didukung}}
\\[0.12cm] $\downarrow$ \\[0.12cm]
\fbox{\parbox{0.78\textwidth}{\centering \textbf{Informasi dan metrik}\\Kriteria, ukuran, dan entitas yang diperlukan}}
\\[0.12cm] $\downarrow$ \\[0.12cm]
\fbox{\parbox{0.78\textwidth}{\centering \textbf{Atribut data}\\Kolom yang dibutuhkan untuk membentuk informasi}}
\\[0.12cm] $\downarrow$ \\[0.12cm]
\fbox{\parbox{0.78\textwidth}{\centering \textbf{Sumber data}\\Tabel, grain, periode, dan relasi antarsumber}}
\\[0.12cm] $\downarrow$ \\[0.12cm]
\fbox{\parbox{0.78\textwidth}{\centering \textbf{Logika pengolahan}\\Join, filter, flag, cleansing, dan transformasi}}
\\[0.12cm] $\downarrow$ \\[0.12cm]
\fbox{\parbox{0.78\textwidth}{\centering \textbf{Validasi dan output}\\Waterfall, duplikasi, konsistensi aturan bisnis, dan hasil akhir}}
\caption{Alur penurunan goal menjadi kebutuhan data}
\label{fig:goal-data}
\end{figure}

Penulis juga mempelajari pola permintaan data. Unit bisnis menyampaikan kebutuhan berdasarkan kriteria. Kriteria tersebut diterjemahkan menjadi query, kemudian hasilnya diperiksa dan dapat melalui proses cleansing sebelum digunakan. Pada tahap ini, penulis belajar bahwa validasi data perlu dilakukan sebelum membahas optimasi query atau visualisasi.

\section{Pengembangan Query SQL dan ETL Campaign}

\subsection{Kebutuhan Pekerjaan}
Salah satu pekerjaan utama adalah mengembangkan query untuk kebutuhan campaign berdasarkan memo atau request dari unit bisnis. Kasus yang dikerjakan mencakup kebutuhan pada produk kartu dan payroll. Query awal mempunyai struktur non-CTE yang panjang dan menggabungkan banyak sumber data.

Skala query berada pada kisaran ratusan baris, menggunakan lebih dari sepuluh sumber tabel, dan menghasilkan lebih dari seratus atribut pada keseluruhan proses. Angka tersebut digunakan untuk menunjukkan kompleksitas pekerjaan, bukan untuk membuka struktur data internal perusahaan.

\subsection{Refactoring Query dengan Common Table Expression}
Penulis merapikan query dengan Common Table Expression atau CTE. Impala mendukung klausa \texttt{WITH} untuk membentuk bagian query yang dapat dibaca sebagai unit terpisah \autocite{impala2026sql}. Pada pekerjaan ini, pemisahan dilakukan berdasarkan fungsi, misalnya pembentukan data dasar, enrichment, flagging, filter, dan output.

Tujuan refactoring bukan sekadar memperpendek query. Struktur CTE membantu penulis menelusuri sumber masalah ketika jumlah baris berubah setelah \textit{join} atau filter. Bagian yang menghasilkan duplikasi dapat diperiksa tanpa membaca seluruh query sebagai satu blok.

Penulis tidak mengubah aturan bisnis secara sepihak. Ketika menemukan kondisi yang ambigu, penulis membandingkan hasil antara query awal dan query hasil refactoring. Perubahan baru dianggap aman jika hasil tetap sesuai dengan tujuan awal atau perbedaannya dapat dijelaskan.

\subsection{Validasi Join, Duplikasi, Flag, dan Waterfall}
Masalah yang paling sering diperiksa adalah duplikasi akibat \textit{join}. Satu identifier dapat muncul lebih dari satu kali pada sumber tertentu. Jika relasi tersebut langsung digabung tanpa memahami kardinalitas, jumlah baris dapat meningkat dan memengaruhi hasil akhir.

Penulis menggunakan beberapa pemeriksaan. Pertama, menghitung jumlah baris dan jumlah identifier unik sebelum dan setelah \textit{join}. Kedua, memeriksa nilai kosong pada atribut yang menjadi syarat. Ketiga, membandingkan hasil filter secara bertahap melalui pendekatan waterfall. Keempat, memeriksa flag yang berasal dari aturan bisnis.

Waterfall membantu menunjukkan titik ketika jumlah data berkurang. Jika pengurangan terlalu besar atau tidak sesuai ekspektasi, penulis kembali ke tahap sebelumnya untuk memeriksa filter. Pendekatan ini lebih mudah ditelusuri dibanding hanya memeriksa hasil akhir.

\subsection{Integrasi ke Talend Data Integration}
Setelah query diuji, pekerjaan dilanjutkan pada Talend Data Integration. Penulis memetakan output query ke struktur dataset yang digunakan pada proses berikutnya. Dataset yang dikerjakan mempunyai puluhan atribut sehingga pemeriksaan tipe data dan urutan mapping menjadi bagian dari pengujian.

Penulis menguji koneksi ke lingkungan Impala, menjalankan job, memeriksa hasil load, dan memastikan komponen yang bergantung pada context dapat menerima parameter yang sesuai. Detail nama koneksi, nama tabel, dan konfigurasi server tidak dituliskan pada laporan karena termasuk informasi internal.

\subsection{Penjadwalan melalui Talend Administration Center}
Job yang bersifat berkala perlu mempunyai mekanisme eksekusi yang konsisten. Talend Administration Center mendukung penjadwalan task menggunakan trigger berbasis waktu atau event \autocite{talend2026}. Pada pekerjaan penulis, skenario yang diuji diarahkan pada eksekusi bulanan sesuai kebutuhan proses campaign.

Pada tahap ini, penulis memahami perbedaan antara job yang berhasil dijalankan secara manual dan job yang siap dijadwalkan. Job yang dijadwalkan perlu mempunyai parameter yang jelas, dependensi yang tersedia, dan hasil yang dapat diperiksa setelah eksekusi.

\section{Analisis Deskriptif Campaign Financing 2026}

\subsection{Tujuan Analisis}
Pekerjaan kedua adalah menyusun analisis deskriptif campaign financing tahun 2026. Tujuannya bukan membangun model prediktif, tetapi membuat data campaign lebih mudah dibaca pada tingkat produk dan campaign.

Analisis menggunakan metrik yang sudah tersedia pada sumber internal, seperti jumlah leads, pesan yang terkirim, conversion, dan outstanding. Definisi conversion mengikuti aturan bisnis pada sumber data. Angka rinci tidak dicantumkan pada laporan karena analisis menggunakan data internal perusahaan.

\subsection{Pembersihan dan Penetapan Scope}
Sebelum membuat perbandingan, penulis memeriksa campaign yang berada di luar fokus analisis dan baris pengiriman yang tidak dapat digunakan. Penulis juga memeriksa definisi beberapa kolom yang dapat menimbulkan salah tafsir.

Langkah ini penting karena rasio dapat terlihat tinggi atau rendah hanya karena scope berbeda. Perbedaan scope tersebut menjadi dasar bagi penulis untuk menyusun asumsi dan pengecualian sebelum membuat ringkasan. Pertanyaan klarifikasi juga disiapkan untuk definisi conversion, periode atribusi, outstanding, duplikasi penerima, dan ketersediaan kelompok pembanding.

\subsection{Agregasi dan Perbandingan}
Analisis dilakukan pada beberapa tingkat. Pertama, penulis menyusun ringkasan keseluruhan untuk melihat ukuran campaign dan hasil secara umum. Kedua, data dikelompokkan per produk. Ketiga, campaign dibandingkan berdasarkan conversion dan ukuran pembiayaan. Keempat, analisis melihat perbedaan mekanisme campaign dan waktu pelaksanaan pada data yang tersedia.

Rasio conversion dihitung dari total conversion dibanding total pesan terkirim pada kelompok yang sama. Penulis tidak menggunakan rata-rata sederhana antarbaris karena volume setiap blast berbeda. Pendekatan agregasi ini menjaga agar campaign kecil tidak mempunyai bobot yang sama dengan campaign besar hanya karena berada pada satu baris.

\subsection{Hasil dan Interpretasi}
Analisis menunjukkan bahwa volume pesan dan hasil campaign tidak selalu bergerak bersama. Terdapat campaign dengan volume besar tetapi conversion relatif rendah. Terdapat juga campaign yang lebih kecil tetapi mempunyai conversion lebih baik. Perbedaan karakter setiap produk membuat penulis tidak menyamakan threshold antarproduk tanpa membaca konteksnya.

Output analisis disusun menjadi ringkasan per produk, daftar campaign yang perlu diperhatikan, serta matriks tindak lanjut. Matriks tersebut digunakan untuk memisahkan campaign yang layak dipertahankan, diperluas secara hati-hati, diperbaiki, atau dievaluasi ulang. Laporan ini tidak mencantumkan nama campaign, angka performa, maupun keputusan internal yang dihasilkan dari matriks tersebut.

\subsection{Batas Interpretasi}
Penulis menemukan beberapa batas yang perlu dijelaskan sebelum hasil digunakan. Pertama, data yang tersedia tidak selalu mempunyai kelompok holdout. Tanpa kelompok pembanding, perubahan conversion tidak dapat langsung disebut sebagai akibat campaign. Kedua, tidak semua data mempunyai waktu kirim yang cukup rinci untuk analisis jam. Ketiga, faktor eksternal dan perubahan perilaku produk tidak dianalisis secara penuh.

Output pekerjaan pada bagian ini diposisikan sebagai analisis deskriptif dan bahan rekomendasi. Penulis menghindari klaim kausal yang tidak didukung desain eksperimen.

\section{Spatial Clustering Data Log Aplikasi Digital}

\subsection{Latar Belakang dan Tujuan}
Pekerjaan ketiga berangkat dari data log lokasi aplikasi digital. Ide awalnya adalah membaca kedekatan lokasi dan pola kunjungan untuk mendukung analisis ekosistem. Dalam pengembangannya, pekerjaan diarahkan pada spatial clustering lebih dahulu karena clustering diperlukan sebelum analisis hubungan yang lebih kompleks.

Social Network Analysis atau SNA ditempatkan sebagai arah lanjutan. Penulis tidak mengklaim SNA sebagai hasil akhir Kerja Praktik karena pengembangan influence score per nasabah belum selesai pada periode laporan.

\subsection{Persiapan Data dan Stay Point}
Data lokasi mempunyai karakter berbeda dari tabel campaign. Satu pengguna dapat menghasilkan banyak ping, sementara koordinat dari perangkat dapat mengalami jitter. Penulis mengubah data menjadi representasi stay point agar aktivitas tidak dibaca hanya dari setiap ping mentah.

Dataset juga dibatasi pada periode dan regional tertentu agar eksperimen dapat dijalankan pada sumber daya komputasi yang tersedia. Batas tersebut membuat hasil tidak boleh dibaca sebagai gambaran seluruh nasabah atau seluruh wilayah.

\subsection{Pemilihan Metode dan Tuning Parameter}
Penulis membandingkan lima metode clustering yang tersedia pada pipeline, yaitu DBSCAN, HDBSCAN, OPTICS, Agglomerative, dan ST-DBSCAN. Kelima metode dijalankan pada data yang sama lalu dinilai menggunakan metrik yang seragam, termasuk DBCV, Silhouette, Davies-Bouldin Index, dan proporsi noise. Kandidat yang tidak memenuhi batas kelayakan operasional dikeluarkan sebelum proses ranking.

Pemilihan metode tidak dilakukan dengan satu nilai metrik saja. Kandidat yang lolos batas kelayakan dibandingkan dengan DBCV sebagai metrik ranking utama dan Silhouette sebagai tie-break. HDBSCAN dipertahankan sebagai metode aktif setelah hasil perbandingan konsisten dengan konfigurasi yang digunakan pada pipeline.

Istilah \textit{brute force} lebih tepat digunakan pada tahap tuning parameter di dalam metode, bukan pada pemilihan lima metode itu sendiri. Setelah HDBSCAN dipilih, penulis menggunakan grid search untuk mencoba seluruh kombinasi parameter yang telah ditetapkan pada ruang pencarian. Setiap kombinasi dijalankan dan dievaluasi, lalu kandidat yang melanggar batas noise atau jumlah cluster dikeluarkan. Pendekatan ini merupakan pencarian exhaustive terhadap grid parameter yang sudah didefinisikan.

Parameter tidak dipilih hanya dari skor tertinggi. Penulis juga memeriksa kestabilan hasil, jumlah cluster, noise, serta apakah perubahan parameter memberi manfaat yang cukup untuk membenarkan perubahan artefak dan interpretasi bisnis. Alur pipeline dan pemilihan metode diringkas pada Gambar \ref{fig:spatial-pipeline}.

\begin{figure}[htbp]
\centering
\small
\setlength{\fboxsep}{6pt}
\fbox{\parbox{0.82\textwidth}{\centering \textbf{Ingest dan feasibility}\\Data log dibaca, kolom disamakan, lalu kelayakan dataset diperiksa}}
\\[0.12cm] $\downarrow$ \\[0.12cm]
\fbox{\parbox{0.82\textwidth}{\centering \textbf{Benchmark lima metode}\\DBSCAN, HDBSCAN, OPTICS, Agglomerative, dan ST-DBSCAN}}
\\[0.12cm] $\downarrow$ \\[0.12cm]
\fbox{\parbox{0.82\textwidth}{\centering \textbf{Evaluasi seragam}\\DBCV, Silhouette, Davies-Bouldin Index, jumlah cluster, dan noise}}
\\[0.12cm] $\downarrow$ \\[0.12cm]
\fbox{\parbox{0.82\textwidth}{\centering \textbf{Constrained method selection}\\Filter kelayakan operasional lalu ranking kandidat}}
\\[0.12cm] $\downarrow$ \\[0.12cm]
\fbox{\parbox{0.82\textwidth}{\centering \textbf{Grid search HDBSCAN}\\Seluruh kombinasi parameter pada grid diuji dan dibandingkan}}
\\[0.12cm] $\downarrow$ \\[0.12cm]
\fbox{\parbox{0.82\textwidth}{\centering \textbf{Output agregat}\\Centroid, reverse geocoding, enrichment, dan visualisasi}}
\caption{Alur pipeline spatial clustering dan pemilihan metode}
\label{fig:spatial-pipeline}
\end{figure}

\subsection{Reverse Geocoding dan Agregasi}
Setelah cluster terbentuk, penulis menghitung centroid untuk mewakili pusat cluster. Reverse geocoding dilakukan pada centroid, bukan pada setiap koordinat pengguna. Cara ini mengurangi kebutuhan untuk mengirim seluruh titik ke layanan geocoding dan menjaga output tetap berada pada tingkat agregat.

Hasil clustering kemudian diperkaya dengan informasi wilayah dan referensi produk yang tersedia. Data disajikan sebagai ringkasan per cluster atau wilayah. Identifier individu tidak ditampilkan pada output laporan.

\subsection{Visualisasi dan Kandidat Ekosistem}
Visualisasi digunakan untuk melihat persebaran cluster, area dengan aktivitas tinggi, jarak terhadap titik referensi, dan komposisi beberapa atribut agregat. Penulis juga menyusun kandidat ekosistem berdasarkan kategori lokasi dari hasil reverse geocoding.

Pada tahap ini, kategori lokasi diperlakukan sebagai interpretasi terhadap centroid. Kategori tersebut bukan atribut individu. Satu area juga dapat mempunyai lebih dari satu fungsi. Karena itu, hasil digunakan sebagai kandidat untuk ditinjau lebih lanjut, bukan sebagai label final terhadap nasabah.

\subsection{Eksperimen Graph dan Arah SNA}
Penulis membangun graph antarcluster berdasarkan kemunculan pengguna yang sama pada lebih dari satu cluster. Graph tersebut digunakan untuk eksplorasi hubungan antararea. Metode community detection juga diuji pada tahap eksplorasi.

Eksperimen tersebut membantu penulis melihat bahwa kedekatan atau kemunculan bersama tidak cukup untuk menyatakan hubungan personal. Co-presence hanya menunjukkan dua entitas muncul pada ruang dan waktu yang berdekatan berdasarkan definisi data. Karena itu, pengembangan SNA memerlukan fitur dan validasi tambahan sebelum dapat menghasilkan influence score yang dapat dipertanggungjawabkan.

\subsection{Batasan Spatial Clustering}
Batasan utama pekerjaan ini adalah kualitas koordinat perangkat, periode observasi yang terbatas, cakupan regional yang terbatas, dan enrichment produk yang belum lengkap. Reverse geocoding juga bergantung pada kualitas data lokasi eksternal.

Keputusan penting pada pekerjaan ini adalah tidak memaksa seluruh ide awal menjadi output. Beberapa tahap yang tidak memberi insight baru dihentikan. SNA juga tidak dinyatakan selesai. Penulis memilih melaporkan batas tersebut daripada membuat kesimpulan yang lebih jauh dari data.

\section{Masalah, Keputusan, dan Pembelajaran Teknis}
Tabel \ref{tab:masalah-keputusan} merangkum beberapa masalah yang berulang pada tiga pekerjaan utama dan keputusan yang diambil penulis.

\begin{table}[htbp]
\caption{Masalah dan keputusan selama pelaksanaan pekerjaan}
\label{tab:masalah-keputusan}
\centering
\begin{tabularx}{\textwidth}{|p{0.27\textwidth}|p{0.33\textwidth}|X|}
\hline
\textbf{Masalah} & \textbf{Risiko} & \textbf{Keputusan} \\
\hline
Join menghasilkan duplikasi & Jumlah target dan agregasi berubah & Periksa kardinalitas, identifier unik, dan waterfall sebelum melanjutkan \\
\hline
Definisi bisnis belum pasti & Query benar secara sintaks tetapi salah secara konteks & Klarifikasi definisi dan dokumentasikan asumsi \\
\hline
Volume campaign berbeda jauh & Rata-rata antarbaris memberi bobot yang tidak seimbang & Gunakan agregasi total pada kelompok yang sama \\
\hline
Tidak ada holdout pada data analisis & Hasil dapat disalahartikan sebagai efek kausal & Batasi kesimpulan pada analisis deskriptif \\
\hline
Koordinat perangkat mengalami jitter & Ping mentah dapat membentuk pola semu & Gunakan stay point, clustering berbasis kepadatan, dan validasi visual \\
\hline
Ide SNA belum mempunyai fitur yang cukup & Influence score berisiko tidak bermakna & Tempatkan SNA sebagai next step dan jangan klaim sebagai hasil selesai \\
\hline
\end{tabularx}
\end{table}

\section{Profesionalisme dan Kerahasiaan}
Kerja Praktik membuat penulis melihat bahwa profesionalisme pada pekerjaan data juga berarti tahu kapan detail tidak boleh dipindahkan keluar lingkungan kerja. Selama pengembangan, penulis menggunakan data dan artefak internal untuk validasi. Untuk laporan akademik, penulis hanya membawa struktur pemikiran, metode, dan pembelajaran yang tidak membuka data perusahaan.

Prinsip yang sama diterapkan ketika menggunakan repository. Repository laporan dibuat terpisah dari repository pekerjaan analitik. Query produksi, file data, kredensial, identifier, dan screenshot internal tidak dimasukkan ke repository laporan. Lampiran juga hanya memuat ringkasan yang telah disanitasi.
